%%%%%%%%%%%%%%%%%%%%%%%%%%%%%%%%%%%%%%%%%%%%%%%%%%%%%%%%%%%%%%%%%%%%%%
% Lớp 12 | Tọa độ của véc-tơ trong không gian | Tọa độ của véc-tơ và các công thức
% Dạng 2H2?2-4: Công thức tọa độ của tích vô hướng và ứng dụng
% Số câu: 48  (sinh tự động bởi tach_cau.py; từ nay thầy sửa trực tiếp file này)
%%%%%%%%%%%%%%%%%%%%%%%%%%%%%%%%%%%%%%%%%%%%%%%%%%%%%%%%%%%%%%%%%%%%%%

%%%%%%%%%%%%%%%% TRẮC NGHIỆM 4 PHƯƠNG ÁN %%%%%%%%%%%%%%%%

% nguồn: Lop_12_Hinh_hoc_Chuong_2
\begin{ex}%[2H2N2-4]%[ma:54cf7fa5]
	Trong không gian $Oxyz$, cho $\vec{u}=(2;-1;-3)$, $\vec{v}=(1;3;-2)$.  Tích vô hướng $\vec{u}\cdot \vec{v}$ bằng
	 \choice
		{$-7$} 
		{\True $5$} 
		 {$-5$}
		{$7$} 
	\loigiai{ 
	Ta có  $\vec{u}\cdot\vec{v}=2\cdot 1+(-1)\cdot 3+(-3)\cdot(-2)=5$.}
\end{ex}

% nguồn: Lop_12_Hinh_hoc_Chuong_2
\begin{ex}%[2H2N2-4]%[ma:fd25a3aa]
	Trong không gian $Oxyz$,  cho hai đường thẳng $d_{1}\colon \dfrac{x}{-1}=\dfrac{y+1}{1}=\dfrac{z-9}{-2}$,\break$d_{2}\colon \dfrac{x+1}{-1}=\dfrac{y-1}{1}=\dfrac{z+5}{1}$. Góc giữa hai đường thẳng $d_1$ và $d_2$ bằng
	 \choice
		{\True  $90^{\circ}$}
		 {$60^{\circ}$}
		{$45^{\circ}$} 
		 {$30^{\circ}$}
	 \loigiai{
		Đường thẳng  $d_{1}$, $d_{2}$ có vectơ chỉ phương lần lượt là $\overrightarrow{u}_1=(-1;1;-2)$, $\overrightarrow{u}_2=(-1;1;1)$.\\
		Ta có 
		$$ 
		\cos (d_{1},d_{2})=\dfrac{\left|\overrightarrow{u}_1\cdot \overrightarrow{u}_2\right|}{|\overrightarrow{u}_1|\cdot|\overrightarrow{u}_2|} 
		 =\dfrac{|(-1)\cdot(-1)+1\cdot1+(-2)\cdot1|}{\sqrt{(-1)^{2}+1^{2}+(-2)^{2}}\cdot \sqrt{(-1)^{2}+1^{2}+1^{2}}}
		=\dfrac{|1+1-2|}{\sqrt{6}\cdot  \sqrt{3}}=0.
		$$ 
		Do đó góc giữa hai  đường thẳng đã cho là $90^{\circ}$.
	} 
\end{ex}

% nguồn: Lop_12_Hinh_hoc_Chuong_2
\begin{ex}%[2H2N2-4]%[ma:2db9c5f6]
	Trong không gian  $Oxyz$, cosin của góc tạo bởi hai vectơ $\overrightarrow{a}=(-1;2;0)$ và $\overrightarrow{b}=(0;-2;1)$ là
	\choice 
		 {$\dfrac{4}{5}$}
		{\True $-\dfrac{4}{5}$} 
		{$\dfrac{4}{25}$} 
		 {$-\dfrac{4}{25}$}
	 \loigiai{
		Ta có $\cos \left(\overrightarrow{a},\overrightarrow{b}\right) =\dfrac{\overrightarrow{a} \cdot \overrightarrow{b}}{\left|\overrightarrow{a}\right|\cdot\left| \overrightarrow{b}\right|}=\dfrac{(-1)  \cdot 0+2 \cdot(-2)+0 \cdot 1}{\sqrt{(-1)^2+2^2+0^2} \cdot \sqrt{0^2+(-2)^2+1^2}}=-\dfrac{4}{5}$.
	 }
\end{ex}

% nguồn: Lop_12_Hinh_hoc_Chuong_2
\begin{ex}%[2H2N2-4]%[ma:62b5e30a]
	Trong không gian  $Oxyz$, cho hai vectơ $\overrightarrow{a}=(-1;3;2)$, $\overrightarrow{b}=(-3;-1;2)$. Tính $\overrightarrow{a} \cdot \overrightarrow{b}$.
	\choice 
		{$10$} 
		{$2$} 
		 {\True $4$}
		{$3$} 
	\loigiai{ 
		Ta có $\overrightarrow{a}  \cdot \overrightarrow{b}=(-1) \cdot(-3)+3 \cdot(-1)+2 \cdot 2=4$.
	} 
\end{ex}

% nguồn: Lop_12_Hinh_hoc_Chuong_2
\begin{ex}%[2H2N2-4]%[ma:7b6a4d93]
	Trong không gian $Oxyz$, cho hai véctơ $\overrightarrow{a}=(1;-2;1)$ và $\overrightarrow{b}=(2;-4;-2)$.  Khi đó $\overrightarrow{a} \cdot \overrightarrow{b}$ bằng
	\choice 
		{\True $8$} 
		{$-8$} 
		{$12$} 
		 {$-12$}
	\loigiai{ 
		Ta có $\overrightarrow{a}  \cdot \overrightarrow{b}=1 \cdot 2+(-2) \cdot(-4)+1 \cdot(-2)=8$.
	 }
\end{ex}

% nguồn: Lop_12_Hinh_hoc_Chuong_2
\begin{ex}%[2H2N2-4]%[ma:dfe0a124]
	Trong không gian với hệ tọa độ  $Oxyz$, cho hai điểm $A(1;-3;1)$, $B(3;0;-2)$. Tính độ dài $AB$.
	\choice 
		 {$26$}
		{$22$} 
		{\True $\sqrt{22}$} 
		 {$\sqrt{26}$}
	 \loigiai{
		Ta có  $\overrightarrow{AB}=(2;3;-3)\Rightarrow AB=\sqrt{2^2+3^2+(-3)^2}=\sqrt{22}$.
	} 
\end{ex}

% nguồn: Lop_12_Hinh_hoc_Chuong_2
\begin{ex}%[2H2N2-4]%[ma:00d985f0]
	Trong không gian $Oxyz$, cho hai vectơ $\overrightarrow{u}=(-3;1;-1)$ và $\overrightarrow{v}=(1;0;5)$. Tích vô  hướng của hai vectơ này bằng
	 \choice
		{\True  $-8$}
		{$8$} 
		{$3$} 
		{$-3$} 
	 \loigiai{
		Tích vô hướng của hai vectơ $\overrightarrow{u}=(-3;1;-1)$ và $\overrightarrow{v}=(1;0;5)$ là 
		$$\overrightarrow{u} \cdot  \overrightarrow{v}=(-3) \cdot 1+1 \cdot 0+(-1) \cdot 5=-8.$$
	} 
\end{ex}

% nguồn: Lop_12_Hinh_hoc_Chuong_2
\begin{ex}%[2H2N2-4]%[ma:33974200]
	Trong không gian $Oxyz$, tích vô hướng của  hai vectơ $\vec{u}=(2; 3;-3)$ và $\vec{v}=(-3; -2; 4)$ bằng
	\choice 
		{$\sqrt{22} \cdot \sqrt{29}$} 
		{$-\sqrt{22}  \cdot \sqrt{29}$}
		{$24 $} 
		{\True $ -24 $} 
	 \loigiai{
		Tích vô  hướng của hai vectơ $\vec{u}$ và $\vec{v}$~là \[\vec{u} \cdot \vec{v} = (2)\cdot(-3)+(3)\cdot(-2)+(-3)\cdot(4)=-6-6-12=-24.\]
	} 
\end{ex}

% nguồn: Lop_12_Hinh_hoc_Chuong_2
\begin{ex}%[2H2N2-4]%[ma:bfd6cbd0]
	Trong không gian $Oxyz$, khoảng  cách giữa hai điểm $I(3; 5; -7)$ và $K(-5; 5; -1)$ bằng
	 \choice
		 {$100$}
		{$20$} 
		 {\True$10$}
		 {$17$}
	\loigiai{Khoảng cách giữa hai điểm $I$ và $K$~là độ dài đoạn thẳng $IK$. Mà \[IK = \sqrt{(-5-3)^2  + (5-5)^2 +\left(-1-(-7)\right)^2}= 10.\]
	} 
\end{ex}

% nguồn: Lop_12_Hinh_hoc_Chuong_2
\begin{ex}%[2H2N2-4]%[ma:6d8e1d71]
	Trong không gian $Oxyz$, cho hai vectơ $\overrightarrow{u}=(-1;3;2)$ và $\overrightarrow{v}=(-3;-1;2)$. Tích vô hướng $\overrightarrow{u}  \cdot \overrightarrow{v}$ bằng
	 \choice
		 {$3$}
		 {$2$}
		{$10$} 
		{\True $4$} 
	\loigiai{ 
		 Ta có $\overrightarrow{u} \cdot \overrightarrow{v}=3-3+4=4$.
	 }
\end{ex}

% nguồn: Lop_12_Hinh_hoc_Chuong_2
\begin{ex}%[2H2H2-4]%[ma:8f347a21]
	 \immini{
	Cho  lập phương $ABCD.A'B'C'D'$ (hình minh họa bên) có độ dài mỗi cạnh bằng $1$. Tính độ dài của vectơ $\overrightarrow{AC}+\overrightarrow{C'D'}$.
	 \choice
		{$\sqrt{3}$} 
		 {$\sqrt{2}$}
		 {\True $1$}
		 {$2\sqrt{2}$}
		} 
		{ 
		\begin{tikzpicture}[scale=0.7,font=\footnotesize,line join=round,line cap=round,>=stealth]
				\coordinate (A) at (0,0); \coordinate (B) at (-2,-1.5); \coordinate (C) at  (3,-1.5); \coordinate (D) at ($(A)+(C)-(B)$);
				\coordinate (A')  at ($(A)+(0,4)$);
				\coordinate (B') at ($(A')+(B)-(A)$); 
				\coordinate  (C') at ($(B')+(C)-(B)$);
				\coordinate (D') at ($(A')+(D)-(A)$); 
				\draw  (D')--(A')--(B')--(C')--(D')--(D)--(C)--(B)--(B')
					 (C)--(C')--(A');
				 \draw[dashed] (D)--(A)--(B)
					 (A')--(A)--(C);
				\foreach  \x/\g in {A/180,B/-90,C/-90,D/0,A'/120,B'/120,C'/90,D'/0} \fill (\x) circle (1pt)
					 ($(\x)+(\g:4mm)$) node{$\x$};
			\end{tikzpicture}
		} 
	 \loigiai{
		Ta  có $A'C'CA$ là hình chữ nhật nên $\overrightarrow{A'C'}=\overrightarrow{AC}$.\\
		Khi đó $\overrightarrow{AC}+\overrightarrow{C'D'}=\overrightarrow{A'C'}+\overrightarrow{C'D'}=\overrightarrow{A'D'}$ nên $\left|  \overrightarrow{AC}+\overrightarrow{C'D'} \right|=\left| \overrightarrow{A'D'} \right|=A'D'=1$.
	 }
\end{ex}

% nguồn: Lop_12_Hinh_hoc_Chuong_2
\begin{ex}%[2H2H2-4]%[ma:75a52f3a]
	 Trong không gian $Oxyz$, cho hai vectơ $\overrightarrow{a}=2\overrightarrow{i} - 3\overrightarrow{k}$ và $\overrightarrow{b} = 4\overrightarrow{i} + \overrightarrow{j} + \overrightarrow{k}$. Tính độ dài của vectơ $\overrightarrow{v}=2\overrightarrow{a}-\overrightarrow{b}$.
	 \choice
		{$\sqrt{74}$} 
		{$3\sqrt{6}$} 
		{\True  $5\sqrt{2}$}
		 {$\sqrt{42}$}
	 \loigiai{
		Ta có $\overrightarrow{a} = (2;0;-3)$ và $\overrightarrow{b} = (4;1;1)$.\\ 
		Khi đó $\overrightarrow{v}=2\overrightarrow{a}-\overrightarrow{b}=2\cdot (2;0;-3)-(4;1;1)=(4;0;-6)-(4;1;1)=(0;-1;-7)$.\\ 
		Vậy độ dài  của vectơ $\overrightarrow{v}$ là $\left|\overrightarrow{v}\right| =\sqrt{0^2+(-1)^2+(-7)^2}=\sqrt{1+49}=\sqrt{50}=5\sqrt{2}$.
	} 
\end{ex}

% nguồn: Lop_12_Hinh_hoc_Chuong_2
\begin{ex}%[2H2H2-4]%[ma:e57072f4]
	Trong không gian với hệ tọa độ $Oxyz$, cho hình chóp $S.ABCD$, đáy $ABCD$ là hình chữ nhật. Biết $A(0;0;0)$, $D(2;0;0)$,  $B(0;4;0)$, $S(0;0;4)$. Gọi $M$ là trung điểm của $SB$ và $G$ là trọng tâm của tam giác $SCD$. Độ dài $MG$ là
	\choice 
		{$MG=\dfrac{\sqrt{6}}{3}$} 
		 {$MG=\dfrac{\sqrt{6}}{2}$}
		{$MG=\dfrac{2\sqrt{3}}{3}$} 
		{\True  $MG=\dfrac{2\sqrt{6}}{3}$}
	 \loigiai{
		\begin{center}
			\begin{tikzpicture}[line join=round,line cap=round,>=stealth,scale=0.65,font=\footnotesize]
					\foreach \x/\y/\n in  {0/0/A,6/0/B,-2/-2/D} \coordinate (\n) at (\x,\y);
					\coordinate (C) at ($(B)+(D)-(A)$); 
					 \coordinate (S) at ($(A)+(0,5)$);
					\coordinate (O) at  ($(B)!1/2!(D)$);
					\coordinate  (M) at ($(S)!1/2!(B)$);
					\coordinate (N)  at ($(C)!1/2!(D)$);
					\coordinate  (G) at ($(S)!2/3!(N)$);
					\coordinate (x) at  ($(A)!3/2!(D)$);
					 \coordinate (y) at ($(A)!5/4!(B)$);
					\coordinate (z) at  ($(A)!4/3!(S)$);
					 \draw (S)--(D)--(C)--(B)--(S)--(C)
						(S)--(N); 
					\draw[dashed]  (S)--(A)--(D)--(B)--(A)--(C)
						 (M)--(G);
					\foreach  \t/\g in {A/180,D/180,C/0,B/45,O/-110,S/135,M/60,G/160}{
						\fill (\t) circle (1pt) node[shift={(\g:8pt)}]{$  \t $};
					 }
					\draw[arrow style] (D)  node[shift={(-70:8pt)}]{$2$}--(x) node[shift={(110:8pt)}]{$x$};
					\draw[arrow style] (B) node[shift={(-70:8pt)}]{$4$}--(y)  node[shift={(90:8pt)}]{$y$};
					\draw[arrow style] (S)  node[shift={(40:8pt)}]{$4$}--(z) node[shift={(0:8pt)}]{$z$};
				\end{tikzpicture}
		\end{center}
		Chọn hệ trục tọa  độ như hình vẽ, ta có $A(0;0;0)$, $D(2;0;0)$, $B(0;4;0)$, $S(0;0;4)$.\\
		$M$ là trung  điểm của $SB\Rightarrow M(0;2;2)$.\\
		Tứ giác $ABCD$ là hình chữ nhật  nên $\heva{
		 &x_A+x_C=x_B+x_D\\
		&y_A+y_C=y_B+y_D\\ 
		 &z_A+z_C=z_B+z_D}
		\Rightarrow\heva{ 
		&x_C=2\\ 
		&y_C=4\\ 
		 &z_C=0}\Rightarrow C(2;4;0)$.\\
		$G$ là trọng tâm của tam giác $SCD\Rightarrow  G\left(\dfrac{4}{3};\dfrac{4}{3};\dfrac{4}{3}\right)$.\\
		Độ dài $MG$  là $MG=\sqrt{\left(\dfrac{4}{3}-0\right)^2+\left(\dfrac{4}{3}-2\right)^2+\left(\dfrac{4}{3}-2\right)^2}=\dfrac{2\sqrt{6}}{3}$.
	} 
\end{ex}

% nguồn: Lop_12_Hinh_hoc_Chuong_2
\begin{ex}%[2H2H2-4]%[ma:c6ec84f6]
	Trong không gian $Oxyz$, cho $\triangle ABC$ có  $\overrightarrow{A B}=(-3;0;4)$, $\overrightarrow{A C}=(5;-2;4)$. Độ dài đường trung tuyến $A M$ là
	\choice 
		{\True $3 \sqrt{2}$} 
		{$5  \sqrt{2}$}
		{$4  \sqrt{2}$}
		 {$2 \sqrt{3}$}
	\loigiai{\immini{Ta  có $$\overrightarrow{A B}+\overrightarrow{A C}=(2;-2;8) \Rightarrow \overrightarrow{A M}=\dfrac{1}{2}\left(\overrightarrow{A B}+\overrightarrow{A C}\right)=(1;-1;4).$$
		Khi đó $A M=\left|\overrightarrow{A M}\right|=\sqrt{1^2+(-1)^2+4^2}=3 \sqrt{2}$.}{\begin{tikzpicture}[scale=.7]
		 \coordinate [label=below:$B$](B) at (0,0);
		\coordinate [label=above:$A$](A) at (2,5); 
		\coordinate [label=below:$C$](C)  at (10,0);
		\coordinate[label=below:$M$] (M) at ($(B)!.5  ! (C)$);
		\foreach  \diem in {A,B,C,M} \fill (\diem) circle (1pt);
		 \draw[smooth](A)--(B)--(C)--(A)--(M);
	\end{tikzpicture}
		} 
	 }
\end{ex}

% nguồn: Lop_12_Hinh_hoc_Chuong_2
\begin{ex}%[2H2H2-4]%[ma:809f6109]
	Trong không gian $Oxyz$ cho hai vectơ $\overrightarrow{u}$, $\overrightarrow{v}$ thoả mãn $\left|\overrightarrow{u}\right|=3$; $\left|\overrightarrow{v}\right|=4$; $\left(\overrightarrow{u};\overrightarrow{v}\right)=60^{\circ}$. Tính độ dài vectơ $\overrightarrow{u}+2  \overrightarrow{v}$.
	 \choice
		{\True  $\sqrt{97}$}
		{$8$} 
		 {$7$}
		{$4  \sqrt{6}$}
	\loigiai{Ta có $\overrightarrow{u} \cdot \overrightarrow{v}=\left|\overrightarrow{u}\right| \cdot\left|\overrightarrow{v}\right| \cos \left(\overrightarrow{u}, \overrightarrow{v}\right)=3\cdot 4\cdot \dfrac{1}{2}=6$. Suy ra 
		$$\left(\left|\overrightarrow{u}+2 \overrightarrow{v}\right|\right)^2=\left|\overrightarrow{u}\right|^2+4 \overrightarrow{u} \cdot \overrightarrow{v}+4\left|\overrightarrow{v}\right|^2=9+24+64=97 \Rightarrow\left|\overrightarrow{u}+2 \overrightarrow{v}\right|=\sqrt{97}.$$ 
	 }
\end{ex}

% nguồn: Lop_12_Hinh_hoc_Chuong_2
\begin{ex}%[2H2H2-4]%[ma:56a81886]
	Trong không gian $Oxyz$, cho hai vec $\overrightarrow{u}=(-2;1;5)$ và $\overrightarrow{v}=(m-2;3;m+1)$, $m$ là tham số. Tìm $m$ để $\overrightarrow{u}$ vuông góc với $\overrightarrow{v}$. 
	\choice 
		{\True  $m=-4$}
		 {$m=4$}
		{$m=-3$} 
		 {$m=3$}
	\loigiai{Ta có $$\overrightarrow{u}\perp\overrightarrow{v} \Leftrightarrow \overrightarrow{u} \cdot \overrightarrow{v}=0 \Leftrightarrow-2 (m-2)+3+5 (m+1)=0  \Leftrightarrow 3 m+12=0 \Leftrightarrow m=-4.$$
	 }
\end{ex}

% nguồn: Lop_12_Hinh_hoc_Chuong_2
\begin{ex}%[2H2H2-4]%[ma:187b596f]
	Trong không gian với hệ tọa độ $Oxyz$, cho ba điểm $A(1;0;2)$, $B(1;1;1)$, $C(2;-1;3)$. Hỏi cosin của góc tạo bởi hai đường thẳng $A B$ và $B C$ bằng bao nhiêu? 
	\choice 
		{$\dfrac{1}{3}$} 
		{$\dfrac{2  \sqrt{2}}{3}$}
		{$\dfrac{\sqrt{3}}{2}$} 
		{$\dfrac{1}{2}$} 
	\loigiai{Ta có $\overrightarrow{B A}=(0;-1;1)$, $\overrightarrow{B  C}=(1;-2;2)$. Suy ra $$ \cos \left(\overrightarrow{B A}, \overrightarrow{B C}\right)=\dfrac{\overrightarrow{B A} \cdot \overrightarrow{B C}}{B A \cdot B C}=\dfrac{4}{3 \sqrt{2}}=\dfrac{2 \sqrt{2}}{3}.$$
		 Vậy $\cos (A B, B C)=\dfrac{2 \sqrt{2}}{3}$.
	 }
\end{ex}

% nguồn: Lop_12_Hinh_hoc_Chuong_2
\begin{ex}%[2H2H2-4]%[ma:4b1cbce9]
	Trong không gian $Oxyz$, cho các vectơ $\overrightarrow{a}=(5;3;-2)$ và $\overrightarrow{b}=(m;-1;m+3)$. Có bao nhiêu giá trị nguyên dương  của $m$ để góc giữa hai vectơ $\overrightarrow{a}$ và $\overrightarrow{b}$ là góc tù?
	 \choice
		 {\True $2$}
		 {$3$}
		 {$1$}
		{$5$} 
	 \loigiai{
		Ta  có $\cos\left(\overrightarrow{a},\overrightarrow{b}\right)=\dfrac{\overrightarrow{a}\cdot\overrightarrow{b}}{\left|\overrightarrow{a}\right|\cdot\left|\overrightarrow{b}\right|}=\dfrac{3m-9}{\sqrt{38}\cdot\sqrt{2m^2+6m+10}}$.\\
		Góc giữa hai vectơ $\overrightarrow{a}$  và $\overrightarrow{b}$ là góc tù khi và chỉ khi $\cos\left(\overrightarrow{a},\overrightarrow{b}\right)<0\Leftrightarrow 3m-9<0\Leftrightarrow m<3$.\\
		Vì $m$ nguyên dương  nên $m\in\left\{1;2\right\}$.\\
		Vậy có $2$ giá trị $m$ thỏa mãn yêu  cầu bài toán.
	 }
\end{ex}

% nguồn: Lop_12_Hinh_hoc_Chuong_2
\begin{ex}%[2H2H2-4]%[ma:7e18c139]
	Trong không gian với hệ trục tọa độ $Oxyz$, biết $|\overrightarrow{u}|=2$, $|\overrightarrow{v}|=1$ và góc giữa  hai vectơ bằng $\dfrac{2 \pi}{3}$. Tìm $k$ để vectơ $\overrightarrow{p}=k \overrightarrow{u}+\overrightarrow{v}$ vuông góc với vectơ $\overrightarrow{q}=\overrightarrow{u}-\overrightarrow{v}$.
	 \choice
		{$k=-\dfrac{2}{5}$} 
		{$k=\dfrac{5}{2}$} 
		{$k=2$} 
		{\True  $k=\dfrac{2}{5}$}
	\loigiai{Ta có $\overrightarrow{p}=k \overrightarrow{u}+\overrightarrow{v}$ vuông góc với vectơ $\overrightarrow{q}=\overrightarrow{u}-\overrightarrow{v}$ khi và  chỉ khi
		\allowdisplaybreaks 
		\begin{eqnarray*}
			\overrightarrow{p} \cdot \overrightarrow{q}=0 &\Leftrightarrow&\left(k \overrightarrow{u}+\overrightarrow{v}\right)\left(\overrightarrow{u}-\overrightarrow{v}\right)=0 \\&\Leftrightarrow& k\left|\overrightarrow{u}\right|^2-\left|\overrightarrow{v}\right|^2+(-k+1)  \overrightarrow{u} \cdot \overrightarrow{v}=0 \\
			&\Leftrightarrow& k\left|\overrightarrow{u}\right|^2-\left|\overrightarrow{v}\right|^2+(-k+1)\left|\overrightarrow{u}\right|  \cdot\left|\overrightarrow{v}\right| \cos \dfrac{2 \pi}{3}=0 \\&\Leftrightarrow& 4 k-1-(-k+1)=0 \Rightarrow k=\dfrac{2}{5} .
		\end{eqnarray*}
	 }
\end{ex}

% nguồn: Lop_12_Hinh_hoc_Chuong_2
\begin{ex}%[2H2H2-4]%[ma:5f498655]
	Trong không gian $Oxyz$, cho  hai vectơ $\overrightarrow{a}=(3;-2;m)$, $\overrightarrow{b}=(2;m;-1)$ với $m$ là tham số nhận giá trị thực. Tìm giá tri của $m$ để hai vectơ $\overrightarrow{a}$ và $\overrightarrow{b}$ vuông góc với nhau.
	\choice 
		 {$m=1$}
		 {\True $m=2$}
		 {$m=-1$}
		 {$m=-2$}
	\loigiai{ 
		Hai vectơ $\overrightarrow{a}$  và $\overrightarrow{b}$ vuông góc với nhau khi và chỉ khi
		$$\overrightarrow{a} \cdot \overrightarrow{b}=0 \Leftrightarrow 3 \cdot 2+(-2) \cdot  m+m \cdot(-1)=0\Leftrightarrow 6-3m=0 \Leftrightarrow m=2.$$
	 }
\end{ex}

% nguồn: Lop_12_Hinh_hoc_Chuong_2
\begin{ex}%[2H2H2-4]%[ma:9ca53015]
	Trong hệ  tọa độ $Oxyz$, cho $\vec{u}=\vec{i}+3\vec{j}+2\vec{k}$ và $\vec{v}=(2;-1;-4)$. Tính $\vec{u}\cdot\vec{v}$.
	 \choice
		{\True  $\vec{u}\cdot \vec{v}=-9$}
		{$\vec{u}\cdot\vec{v}=9$} 
		{$\vec{u}\cdot\vec{v}=(2;-3;1)$} 
		{$\vec{u}\cdot\vec{v}=\sqrt{14}$} 
	 \loigiai{
		Từ $\vec{u}=\vec{i}+3\vec{j}+2\vec{k}\Rightarrow  \vec{u}=(1;3;2)$.\\
		Do đó,  $\vec{u}\cdot\vec{v}=1\cdot 2+3\cdot(-1)+2\cdot(-4)=-9$.
	 }
\end{ex}

% nguồn: Lop_12_Hinh_hoc_Chuong_2
\begin{ex}%[2H2H2-4]%[ma:4c502aff]
	Cho hai vectơ $\overrightarrow{a}$ và $\overrightarrow{b}$ tạo với nhau một góc $60^{\circ}$ và  $\left|\overrightarrow{a}\right|=2$; $\left|\overrightarrow{b}\right|=4$. Khi đó $\left|\overrightarrow{a}+\overrightarrow{b}\right|$ bằng
	\choice 
		{$2 \sqrt{5}$} 
		{\True  $2 \sqrt{7}$}
		 {$2$}
		{$\sqrt{8  \sqrt{3}+20}$}
	\loigiai{ 
		Ta có 
		 \allowdisplaybreaks
		\begin{eqnarray*}
			\left|\overrightarrow{a}+\overrightarrow{b}\right|^2&=&\left(\overrightarrow{a}+\overrightarrow{b}\right)^2=\left(\overrightarrow{a}\right)^2+2 \overrightarrow{a} \overrightarrow{b}+\left(\overrightarrow{b}\right)^2\\&=&\left|\overrightarrow{a}\right|^2+2\left|\overrightarrow{a}\right|\cdot \left|\overrightarrow{b}\right| \cos \left(\overrightarrow{a}, \overrightarrow{b}\right)+\left|\overrightarrow{b}\right|^2\\&=&4+2  \cdot 2 \cdot 4 \cdot \dfrac{1}{2}+16=28\\ \Rightarrow\left|\overrightarrow{a}+\overrightarrow{b}\right|&=&\sqrt{28}=2 \sqrt{7}.
		\end{eqnarray*}
	} 
\end{ex}

% nguồn: Lop_12_Hinh_hoc_Chuong_2
\begin{ex}%[2H2V2-4]%[ma:c62c9194]
	Trong không gian $Oxyz$, cho ba điểm $A(1;1;1)$, $B(2;3;2)$, $C(3;-1;3)$. Tìm tọa  độ điểm $D$ sao cho bốn đỉnh $A$, $B$, $C$, $D$ lập thành một hình chữ nhật.
	\choice 
		 {$D(4;3;4)$}
		{$D(4;-1;4)$} 
		{$D(2;-3;2)$} 
		 {\True $D(4;1;4)$}
	 \loigiai{
		Ta có $\overrightarrow{AB} = (1;2;1)$, $\overrightarrow{AC}=(2;-2;2)$.\\ 
		Do  $\overrightarrow{AB}\cdot \overrightarrow{AC} = 1\cdot 2 + 2\cdot (-2) + 1\cdot 2 = 0$ nên $AB \perp AC$.\\
		Điều này cho thấy $A$, $B$, $C$ không thẳng hàng và hình chữ nhật tạo ra phải  là $ABDC$.\\
		Gọi $D(x;y;z)$, ta  có $\overrightarrow{CD} = (x-3; y+1; z-3)$.\\
		Tứ giác $ABDC$ là hình chữ nhật khi và  chỉ khi
		$$\overrightarrow{AB} = \overrightarrow{CD}\Leftrightarrow 
		\heva{& x-3=1 \\& y+1=2 \\& z-3=1 } \Leftrightarrow \heva{& x=4 \\& y=1 \\&  z=4} \Rightarrow D(4;1;4).$$
		 Vậy $D(4;1;4)$.
	 }
\end{ex}

% nguồn: Lop_12_Hinh_hoc_Chuong_2
\begin{ex}%[2H2V2-4]%[ma:ae6df059]
	Trong không gian  $Oxyz$, cho ba điểm $A(1;4;-5)$, $B(2;3;-6)$ và $C(4;4;-5)$. Tìm tọa độ trực tâm $H$ của tam giác $ABC$.
	 \choice
		{$H\left(\dfrac{5}{7};4;-5\right)$} 
		 {$H(1;4;-5)$}
		{\True $H(2;3;-6)$} 
		{$H\left(\dfrac{7}{3};\dfrac{11}{3};-\dfrac{16}{3}\right)$} 
	\loigiai{Gọi $H(x;y;z)$ là trực  tâm của $\triangle A B C$.\\
		Ta có $\overrightarrow{A B}=(1;-1;-1)$, $\overrightarrow{A C}=(3;0;0)$, suy ra $\overrightarrow{n}=\left[\overrightarrow{A B}, \overrightarrow{A C}\right]=(0;-3;3)$ là vectơ  pháp tuyến của mặt phẳng $(A B C)$. Suy ra $(A B C)\colon y-z-9=0$.\\
		Vì $H(x;y;z)$ là trực tâm của tam giác $A B C$  nên
		 \allowdisplaybreaks
		\begin{eqnarray*}
			\heva{&\overrightarrow{C H} \cdot \overrightarrow{A B}=0 \\& \overrightarrow{B  H} \cdot \overrightarrow{A C}=0 \\& H \in(A B C)} &\Leftrightarrow&\heva{&(x-4) \cdot 1+(y-4) \cdot(-1)+(z+5) \cdot(-1)=0 \\& (x-2) \cdot 3+(y-3) \cdot 0+(z+6) \cdot 0=0 \\& y-z-9=0} \\&\Leftrightarrow&\heva{&x-y-z=5 \\& x=2 \\& y-z=9} \Leftrightarrow\heva{&x=2 \\& y=3 \\& z=-6.}
		\end{eqnarray*}
		Vậy $H(2;3;-6)$. 
	} 
\end{ex}

% nguồn: Lop_12_Hinh_hoc_Chuong_2
\begin{ex}%[2H2V2-4]%[ma:9c419868]
	Trong không gian với hệ trục tọa độ $Oxyz$, cho hai điểm $A(1;-2;3)$, $B(-2;2;2)$. Gọi  $I(a, b, c)$ là tâm đường tròn ngoại tiếp tam giác $OAB$. Tính $T=a^2+b^2+c^2$.
	 \choice
		{\True $T=\dfrac{13}{2}$} 
		 {$T=6$}
		 {$T=2$}
		{$T=\dfrac{29}{4}$} 
	\loigiai{Theo giả thiết  ta có
		 \allowdisplaybreaks
		\begin{eqnarray*}
			\heva{&I O=I A \\& I O=I B \\& \left[\overrightarrow{O A};\overrightarrow{O B}\right] \cdot \overrightarrow{O I}=0} &\Leftrightarrow&\heva{&a^2+b^2+c^2=(1-a)^2+(-2-b)^2+(3-c)^2 \\& a^2+b^2+c^2=(-2-a)^2+(2-b)^2+(2-c)^2 \\& -10 a-8 b-2  c=0}\\
			&\Leftrightarrow&\heva{&-2 a+4 b-6 c=-14 \\& 4 a-4 b-4  c=-12 \\& -10 a-8 b-2 c=0} \Leftrightarrow\heva{&a=-\dfrac{1}{2} \\& b=0 \\& c=\dfrac{5}{2}.}
		\end{eqnarray*}
		Vậy $T=a^2+b^2+c^2=\dfrac{1}{4}+\dfrac{25}{4}=\dfrac{13}{2}$. 
	 }
\end{ex}


%%%%%%%%%%%%%%%% ĐÚNG/SAI %%%%%%%%%%%%%%%%

% nguồn: Lop_12_Hinh_hoc_Chuong_2
\begin{ex}%[2H2H2-4]%[ma:431fbe09]
	Cho hai vectơ $\vec{u}=\left(0;2;3\right)$  và $\vec{v}=\left(m-1;2m;3\right)$.
	\choiceTF 
		{\True  $\left|\vec{u}\right|=\sqrt{13}$}
		 {$\left|\vec{u}\left|=\right|\vec{v}\right|\Leftrightarrow m=-\dfrac{3}{5}$}
		{\True $\vec{u}=\vec{v}\Leftrightarrow  m=1$}
		{$\vec{u}\perp \vec{v}\Leftrightarrow m=\dfrac{9}{4}$} 
	 \loigiai{
		\begin{itemchoice}
			\itemch $\left|\vec{u}\right|=\sqrt{0^2+2^2+3^2}=\sqrt{13}$. 
			\itemch $\left|\vec{u}\left|=\right|\vec{v}\right|\Leftrightarrow \sqrt{13}=\sqrt{\left(m-1\right)^2+4m^2+9}  \Leftrightarrow 5m^2-2m-3=0\Leftrightarrow m=1$ hoặc $m=-\dfrac{3}{5}$.
			\itemch Khi $m=1$ thì $\vec{v}=\left(0;2;3\right)$. Suy ra  $\vec{u}=\vec{v}$.
			\itemch  $\vec{u}\perp \vec{u}\Leftrightarrow 4m+9=0\Leftrightarrow m=-\dfrac{9}{4}$.
		\end{itemchoice}
	} 
\end{ex}

% nguồn: Lop_12_Hinh_hoc_Chuong_2
\begin{ex}%[2H2H2-4]%[ma:f684e990]
	Trong không gian  $Oxyz$, cho ba điểm $A(-1;2;3)$, $B(3;-4;1)$ và $C(2;-5;1)$.
	 \choiceTF
		{\True  Tọa độ của hình chiếu vuông góc của điểm $A$ trên mặt phẳng $Oxy$ là $(-1;2;0)$}
		{$\overrightarrow{AB}=(2;-3;-1)$} 
		 {\True Đường thẳng $BC$ song song với mặt phẳng $(Oxy)$}
		{Gọi $M$ là điểm trên mặt phẳng $(Oxy)$ sao cho ba điểm $A$, $B$, $M$ thẳng hàng.  Khi đó $CM=\sqrt{15}$}
	 \loigiai{
		\begin{itemchoice}
			\itemch {\bf Đúng.}\\ Vì hình chiếu vuông góc của điểm $A(-1;2;3)$ trên mặt phẳng $Oxy$ là $(-1;2;0)$. 
			\itemch {\bf Sai.}\\ Vì  $\overrightarrow{AB}=(4;-6;-2)$.
			\itemch {\bf Đúng.}\\ $\overrightarrow{BC}=(-1;-1;0)$ và mặt phẳng  $(Oxy)$ có một véc-tơ pháp tuyến là $\overrightarrow{k}=(0;0;1)$.\\
			Vì $\overrightarrow{BC}$ vuông góc với  $\overrightarrow{k}$ và $B\notin (Oxy)$ nên $BC$ song song với $(Oxy)$.
			\itemch {\bf Sai.}\\ Vì gọi $M\left(x_0;y_0;0\right)$ thuộc  mặt phẳng $(Oxy)$.\\
			Do ba điểm $A$, $B$, $M$ thẳng hàng  nên hai véc-tơ $\overrightarrow{AM}$, $\overrightarrow{AB}$ cùng phương, với $\overrightarrow{AM}=\left(x_0+1;y_0-2;-3\right)$, $\overrightarrow{AB}=(4;-6;-2)$.\\
			Suy ra $\dfrac{x_0+1}{4}=\dfrac{y_0-2}{-6}=\dfrac{-3}{-2}\Rightarrow\heva{&x_0=5\\&y_0=-7}\Rightarrow M(5;-7;0)$.\\ 
			Khi đó  $CM=\sqrt{(5-2)^2+(-7+5)^2+(0-1)^2}=\sqrt{14}$.
		\end{itemchoice}
	 }
\end{ex}

% nguồn: Lop_12_Hinh_hoc_Chuong_2
\begin{ex}%[2H2H2-4]%[ma:b9bee03e]
	Trong không gian $Oxyz$, cho hai véctơ $\overrightarrow{a}=(-2;1;-3)$,  $\overrightarrow{b}=(-1;-3;2)$ và điểm $A(4;6;-3)$.
	\choiceTF 
		{Tọa độ vectơ  $\overrightarrow{a}-2\overrightarrow{b}=(0;1;-1)$}
		{\True Tọa  độ điểm $B(2;7;-6)$ thì $\overrightarrow{a}=\overrightarrow{AB}$}
		{Hai vectơ $\overrightarrow{a}$ và $\overrightarrow{b}$ cùng phương hướng} 
		{\True Góc giữa vectơ $\overrightarrow{a}$ và $\overrightarrow{b}$ bằng $120^\circ$} 
	 \loigiai{
		\begin{itemchoice}
			\itemch Ta có $\overrightarrow{a}-2\overrightarrow{b}=(0;7;-7)$ 
			\itemch Ta có  $\overrightarrow{a}=\overrightarrow{AB}\Leftrightarrow\heva{&-2=x_B-4\\&1=y_B-6\\&-3=z_B+3}\Leftrightarrow\heva{&x_B=2\\&y_B=7\\&z_B=-6}\Rightarrow B(2;7;-6)$
			\itemch Ta có $\dfrac{-2}{-1}\neq \dfrac{1}{-3}\neq \dfrac{-3}{2}$ 
			\itemch Ta có $\cos  \left(\overrightarrow{a},\overrightarrow{b}\right)=\dfrac{-2\cdot(-1)+(-3)\cdot 1+(-3)\cdot 2}{\sqrt{2^{2}+1^{2}+(-3)^{2}}\cdot \sqrt{1^{2}+(-3)^{2}+2^{2}}}=-\dfrac{1}{2}\Rightarrow\left(\overrightarrow{a},\overrightarrow{b}\right)=120^\circ$
		\end{itemchoice}
	} 
\end{ex}

% nguồn: Lop_12_Hinh_hoc_Chuong_2
\begin{ex}%[2H2H2-4]%[ma:547f669e]
	Trong không gian $Oxyz$, cho ba điểm $A(3; 5;-1)$, $B(7;x;1)$,  $C(9;2;y)$. Khi đó
	 \choiceTF
		{\True Ba điểm $A$, $B$, $C$ thẳng hàng  thì $x+y=5$}
		{Điểm $G\left(\dfrac{19}{3};\dfrac{8}{3};3\right)$ là trọng  tâm tam giác $ABC$ thì $x=1$, $y=3$}
		{\True  Tích vô hướng $\overrightarrow{AB} \cdot \overrightarrow{AC}=-3x+2y+41$}
		{Tam giác $ABC$ vuông tại $A$ thì $x=13$, $y=-1$} 
	\loigiai{ 
		\begin{itemchoice}
			\itemch Ta có $\overrightarrow{AB}=(4;x-5;2)$, $\overrightarrow{AC}=(6;-3;y+1)$.\\ 
			Để ba điểm $A$, $B$, $C$ thẳng hàng thì  $$\overrightarrow{AB}=k\cdot \overrightarrow{AC} \Leftrightarrow\heva{&4=k\cdot 6 \\& x-5=k(-3) \\&2=k(y+1)} \Leftrightarrow\heva{&k=\dfrac{2}{3} \\& x=3 \\&y=2.}$$
			 Vậy $x+y=5$.
			\itemch Nếu $G\left(\dfrac{19}{3};\dfrac{8}{3};3\right)$ là trọng tâm tam giác $ABC$ thì $\heva{&\dfrac{3+7+9}{3}=\dfrac{19}{3}  \\&\dfrac{5+x+2}{3}=\dfrac{8}{3} \\& \dfrac{-1+1+y}{3}=3}\Leftrightarrow\heva{&x=1 \\&y=9.}$
			\itemch Ta  có $\overrightarrow{AB} \cdot \overrightarrow{AC}=24+(x-5)(-3)+2(y+1)=-3x+2y+41$.
			\itemch Tam giác $ABC$ vuông tại $A$  khi và chỉ khi đó $\overrightarrow{AB} \cdot \overrightarrow{AC}=-3x+2y+41=0$.\\
			Lúc đó cặp $x=1$, $y=19$ cũng thỏa  mãn điều kiện.
		\end{itemchoice}
	 }
\end{ex}

% nguồn: Lop_12_Hinh_hoc_Chuong_2
\begin{ex}%[2H2H2-4]%[ma:ccee839a]
	Cho ba  vectơ $\vec{a}=\left(-1;1;0\right)$, $\vec{b}=\left(1;1;0\right)$ và $\vec{c}=\left(1;1;1\right)$.
	\choiceTF 
		{$\left|\vec{a}\right|=2$} 
		{\True  $\left|\vec{c}\right|=\sqrt{3}$}
		 {$\cos\left(\vec{a},\vec{c}\right)=\dfrac{2}{\sqrt{5}}$}
		{$\vec{b}\perp  \vec{c}$}
	\loigiai{ 
		\begin{itemchoice}
			\itemch $\left|\vec{a}\right|=\sqrt{\left(-1\right)^2+1^2}=\sqrt{2}$. 
			\itemch  $\left|\vec{c}\right|=\sqrt{1^2+1^2+1^2}=\sqrt{3}$.
			 \itemch $\cos\left(\vec{a},\vec{c}\right)=\dfrac{\vec{a}\cdot \vec{c}}{\left|\vec{a}\right|\cdot\left|\vec{c}\right|}=0$.
			\itemch $\vec{b}\cdot\vec{c}=2$ suy ra $\vec{b}$ không vuông  góc với $\vec{c}$.
		\end{itemchoice}
	} 
\end{ex}

% nguồn: Lop_12_Hinh_hoc_Chuong_2
\begin{ex}%[2H2H2-4]%[ma:d5a0c1f1]
	 Trong không gian $Oxyz$, cho ba điểm $A\left(2;-1;1\right)$, $B\left(-1;3;-1\right)$, $C\left(5;-3;4\right)$. Khi đó
	 \choiceTF
		{\True Tích vô hướng của hai vectơ $\overrightarrow{AB}$ và  $\overrightarrow{BC}$ bằng $-52$}
		{Góc $\widehat{ABC}$ là góc  tù}
		{\True  Côsin góc giữa hai vectơ $\overrightarrow{AB}$, $\overrightarrow{AC}$ bằng $\dfrac{-23}{\sqrt{638}}$}
		{\True Điểm $D\left(1;2;x\right)$ với $\triangle  ABD$ vuông tại $B$ thì giá trị $x=-6$}
	 \loigiai{
		\begin{itemchoice}
			 \itemch
			Ta có $\overrightarrow{AB}=\left(-3;4;-2\right) $; $ \overrightarrow{BC}=\left(6;-6;5\right)$.\\ 
			Suy ra $ \overrightarrow{AB}\cdot\overrightarrow{BC} =\left(-3\right)\cdot 6+4\cdot\left(-6\right)+\left(-2\right)\cdot 5=-52$. 
			\itemch 
			Ta có  $\overrightarrow{BA}\cdot\overrightarrow{BC}=52> 0$ nên góc $\widehat{ABC}$ là góc nhọn.
			\itemch 
			Ta có $\overrightarrow{AB}=\left(-3;4;-2\right)$;  $ \overrightarrow{AC}=\left(3;-2;3\right)$. Suy ra
			\[\cos \left(\overrightarrow{AB}, \overrightarrow{AC}\right)=\dfrac{\overrightarrow{AB}.\; \overrightarrow{AC}}{\left|\overrightarrow{AB}\right|\cdot\left|\overrightarrow{AC}\right|}= \dfrac{-9-8-6}{\sqrt{\left(-3\right)^2+4^2+\left(-2\right)^2}\cdot\sqrt{3^2+\left(-2\right)^2+3^2}}=\dfrac{-23}{\sqrt{638}}.\] 
			 \itemch
			Ta  có $\overrightarrow{AB}=\left(-3;4;-2\right)$; $\overrightarrow{BD}=\left(2;-1;x+1\right)$.\\
			Tam  giác $ABD$ vuông tại $B$ khi và chỉ khi
			\[\overrightarrow{AB}\perp \overrightarrow{BD}\Leftrightarrow  \overrightarrow{AB}\cdot\overrightarrow{BD}=0\Leftrightarrow-6-4-2\left(x+1\right)=0\Leftrightarrow x=-6.\]
		\end{itemchoice}
	 }
\end{ex}

% nguồn: Lop_12_Hinh_hoc_Chuong_2
\begin{ex}%[2H2H2-4]%[ma:aa909a49]
	Cho tam giác $ABC$ có $A\left(1;2;0\right)$,  $B\left(0;1;1\right)$, $C\left(2;1;0\right)$.
	 \choiceTF
		{\True Tam giác $ABC$  vuông tại $A$}
		{Chu  vi tam giác là $\sqrt{7}+\sqrt{3}+\sqrt{2}$}
		{Diện tích tam giác $ABC$ là  $\sqrt{6}$}
		{\True Tâm đường tròn ngoại  tiếp tam giác $ABC$ là $I\left(1;1;\dfrac{1}{2} \right)$}
	 \loigiai{
		Ta có $\overrightarrow{AB}=\left(-1;-1;1\right)\Rightarrow AB=\sqrt{3}$, $\overrightarrow{AC}=\left(1;-1;0\right)\Rightarrow AC=\sqrt{2}$,\\ $\overrightarrow{BC}=\left(2;0;-1\right)\Rightarrow  BC=\sqrt{5}$.
		\begin{itemchoice}
			\itemch $\overrightarrow{AB}\cdot \overrightarrow{AC}=0$  do đó $AB\perp AC$, tam giác $ABC$ vuông tại $A$.
			\itemch Chu vi của tam giác  là $AB+AC+BC=\sqrt{3}+\sqrt{2}+\sqrt{5}$.
			\itemch Diện tích $\triangle ABC$ là $S=\dfrac{1}{2} \cdot  AB\cdot AC=\dfrac{\sqrt{6}}{2}$
			\itemch Tâm  đường tròn ngoại tiếp là trung điểm của $BC$ có tọa độ $I\left(1;1;\dfrac{1}{2} \right)$.
		\end{itemchoice}
	} 
\end{ex}

% nguồn: Lop_12_Hinh_hoc_Chuong_2
\begin{ex}%[2H2V2-4]%[ma:3b7fb600]
	\immini[thm]{Trong không gian với hệ trục tọa độ $Oxyz$, cho hình lăng trụ $OAB.O'A'B'$. Biết $O(0;0;0)$, $A(2;0;0)$,  $B(0;1;0)$, $O'(0;0;3)$.
	 \choiceTF
		{\True Đường thẳng $AO'$ có một vectơ chỉ phương là  $\overrightarrow{a}=(2;0;-3)$}
		{Góc gữa hai đường thẳng $O'A'$ và $AB$ bằng $56^\circ28'$ (làm  tròn kết quả đến hàng phút)}
		 {\True Mặt phẳng $(ABO')$ có một vectơ pháp tuyến là $\overrightarrow{n}=(3;6;2)$}
		{\True Trong tất các mặt cầu tiếp xúc với hai đường thẳng $AB$  và $OO'$ thì mặt cầu có bán kính $R=\dfrac{\sqrt{5}}{5}$ là mặt cầu có bán kính nhỏ nhất trong các mặt cầu nói trên}}
		{\begin{tikzpicture}[scale=0.7,font=\footnotesize,line join=round,line cap=round,>=stealth]
		\path 
			(0,0) coordinate  (O)
			(-1.5,-2) coordinate  (A)
			(3,0)  coordinate (B)
			($(O)+(0,5)$)  coordinate (O')
			 ($(O')+(A)-(O)$) coordinate (A')
			($(O')+(B)-(O)$) coordinate  (B');
		\draw  [->]
			(O')--($(O')+(0,1)$); 
		 \draw [->]
			 (B)--($(B)+(1,0)$);
		\draw [->] 
			 (A)--($(A)+(-0.75,-1)$);
		 \draw (A)--(A')
			 (B)--(B')
			 (A)--(B)
			(A')--(B') 
			 (B')--(O')
			(O')--(A'); 
		 \draw[dashed] (A)--(O)--(B)
			(O)--(O'); 
		\foreach \p/\q in {A/-180,B/-90,O/-45,A'/90,O'/135} \fill (\p)  circle (1pt)
			($(\p)+(\q:3.5mm)$) node{$\p$}; 
		\foreach \p/\q in {B'/90}  \fill (\p) circle (1pt)
			($(\p)+(\q:3.5mm)$)  node{$\p$};
	\end{tikzpicture}
		 }
	\loigiai{ 
	\begin{itemchoice}
		\itemch Ta có  $\overrightarrow{AO'}=(2;0;-3)$.
		\itemch Ta có  $\left|\overrightarrow{AB}\right|=\sqrt{5}$, $\left|\overrightarrow{O'A'}\right|=2$. \\
		Khi đó $\cos\alpha=\dfrac{\left|\overrightarrow{O'A'}\cdot\overrightarrow{AB}\right|}{\left|\overrightarrow{O'A'}\right|\cdot\left|\overrightarrow{AB}\right|}=\dfrac{2\sqrt{5}}{5}$.  \\
		Suy  ra $\alpha\approx 26{,}565\approx 26^\circ33'$.
		\itemch Ta có  $\overrightarrow{AO'}=(-2;0;3)$, $\overrightarrow{AB}=(-2;1;0)$. \\
		Gọi $\overrightarrow{n}$ là $1$ vectơ pháp tuyến của mặt  phẳng cần tìm.\\
		Vì $\heva{& \overrightarrow{n}\perp \overrightarrow{AO'}\\ & \overrightarrow{n}\perp  \overrightarrow{AB}}\Rightarrow \overrightarrow{n}=\left[\overrightarrow{AB},\overrightarrow{AO'}\right]=(3;6;2)$.
		\itemch Mặt cầu có bán kính nhỏ nhất tiếp xúc với hai đường thẳng  có tâm là trung điểm đoạn vuông góc chung. \\
		 Có $\overrightarrow{OO'}=(0;0;3)$, $\overrightarrow{AB}=(-2;1;0)$ \\
		Suy ra Phương trình của  hai đường thẳng là
		\[OO'\colon\heva{& x=0  \\ & y=0 \\ & z=3t};\quad AB\colon\heva{& x=2-2t' \\ & y=t' \\ & z=0.}\]
		Gọi $M\in OO'$; $M(0;0;3t)$, $N\in AB$; $N(2-2t';t';0)$ sao cho $MN$ là đoạn  vuông góc chung giữa $OO'$ và $AB$. \\
		Khi  đó $\overrightarrow{MN}=(2-2t';t';-3t)$. \\
		Do $MN$ là đoạn vuông góc chung \[\Leftrightarrow \heva{& \overrightarrow{MN}\cdot\overrightarrow{OO'}=0 \\ & \overrightarrow{MN}\cdot\overrightarrow{AB}=0} \Leftrightarrow \heva{& -9t=0 \\ & -2(2-2t')+t'=0} \Leftrightarrow \heva{& t=0 \\ & t'=\dfrac{4}{5}} \Rightarrow \heva{&  M=(0;0;0) \\ & N=\left(\dfrac{2}{5};\dfrac{4}{5};0\right).}\]
		Suy ra bán kính nhỏ  nhất mặt cầu thỏa mãn là
		 \[\dfrac{\left|\overrightarrow{MN}\right|}{2}=\dfrac{\sqrt{\left(\dfrac{2}{5}\right)^2+\left(\dfrac{4}{5}\right)^2}}{2}=\dfrac{\sqrt{5}}{5}.\]
	\end{itemchoice}
	 }
\end{ex}

% nguồn: Lop_12_Hinh_hoc_Chuong_2
\begin{ex}%[2H2V2-4]%[ma:098e3093]
	Trong không gian $Oxyz$, cho ba điểm  $A(0;-2;1)$, $B(-2;-2;-1)$, $C(3;1;-2)$. Xét tính đúng sai của các mệnh đề sau
	\choiceTF 
		{\True Hình chiếu của điểm  $C$ lên trục $Oz$ là $C'(0;0;-2)$}
		{Trọng tâm của tam giác  $ABC$ là $G \left(\dfrac{1}{3};1;-\dfrac{2}{3} \right)$}
		{\True Tam  giác $ABC$ vuông tại $A$}
		{Nếu $M(a;b;c)$  là điểm thuộc mặt phẳng $(Oxy)$ sao cho ba điểm $B$, $C$, $M$ thẳng hàng thì $a+b=12$}
	\loigiai{ 
		\begin{itemchoice}
			\itemch 
			Hình chiếu của điểm $C(3;1;-2)$  lên trục $Oz$ là $C'(0;0;-2)$.
			 \itemch
			$G(a;b;c)$ là  trọng tâm của tam giác $ABC$, suy ra $\heva{
			 &x_G=\dfrac{0+(-2)+3}{3}=\dfrac{1}{3}\\&
			 y_G=\dfrac{(-2)+(-2)+1}{3}=-1 \\&
			z_G=\dfrac{1+(-1)+(-2)}{3}=-\dfrac{2}{3}}$.\\ 
			 Vậy $G\left( \dfrac{1}{3};-1;-\dfrac{2}{3}\right)$.
			 \itemch
			Ta có $\overrightarrow{AB}(-2;0;-2)$, $\overrightarrow{AC}(3;3;-3)$. Ta thấy  $\overrightarrow{AB} \cdot \overrightarrow{AC} =0$ nên $AB \perp AC$.\\
			Suy ra tam giác $ABC$ vuông  tại $A$.
			 \itemch
			Ta có $\overrightarrow{BC}(5;3;-1)$, $\overrightarrow{BM}(a+2;b+2;1)$. Do ba điểm $B,C,M$ thẳng hàng nên $\overrightarrow{BC}$, $\overrightarrow{BM}$ cùng phương nên suy ra  $\dfrac{a+2}{5}=\dfrac{b+2}{3}=\dfrac{1}{-1} \Rightarrow \heva{a=-7 &\\
			b=-5} \Rightarrow a+b=-12$. 
		\end{itemchoice}
	 }
\end{ex}


%%%%%%%%%%%%%%%% TRẢ LỜI NGẮN %%%%%%%%%%%%%%%%

% nguồn: Lop_12_Hinh_hoc_Chuong_2
\begin{bt}%[2H2N2-4]%[ma:f3d38592]
	Trong không gian với hệ tọa~độ $Oxyz$, cho hai vectơ $\vec{a}=(1; 1;-3), $ $\vec{b}=(-5;-1; 0)$.  Tích vô hướng của hai vectơ $\vec{a},$ $ \vec{b}$ có giá trị bằng bao nhiêu?
	 \shortans{-6}
	\loigiai{Ta có $$\vec{a} \cdot \vec{b}=1 \cdot (-5)+1 \cdot  (-1)+(-3) \cdot 0=-6.$$
	 }
\end{bt}

% nguồn: Lop_12_Hinh_hoc_Chuong_2
\begin{bt}%[2H2N2-4]%[ma:1aead6ac]
	Trong  không gian với hệ tọa~độ $Oxyz$, cho hai vectơ $\vec{a}=(-1; 3; 2), $ $\vec{b}=(x; 0; 1)$. Giá trị của $x$ để $\vec{a}\perp \vec{b}$~là bao nhiêu?
	\shortans{2} 
	 \loigiai{Ta có
		\allowdisplaybreaks 
		\begin{eqnarray*}
			\vec{a} \cdot \vec{b}&=&(-1) \cdot x+3 \cdot 0+2 \cdot 1=-x+2,\\  \vec{a} \cdot \vec{b}&=&0 \Leftrightarrow x=2.
		\end{eqnarray*}
	 }
\end{bt}

% nguồn: Lop_12_Hinh_hoc_Chuong_2
\begin{bt}%[2H2N2-4]%[ma:470f2975]
	Trong không gian với hệ tọa~độ $Oxyz$, độ dài của vectơ $\vec{a}=(2; 3; 6)$ bằng bao nhiêu? 
	\shortans{7} 
	 \loigiai{
		 Ta có $$\left|\vec{a}\right|=\sqrt{2^2+3^2+6^2}=7.$$
	} 
\end{bt}

% nguồn: Lop_12_Hinh_hoc_Chuong_2
\begin{bt}%[2H2H2-4]%[ma:95cce52a]
	Trong không gian với  hệ tọa~độ $Oxyz$, cho hai vectơ $\vec{a}=(0;-1; 0),$ $ \vec{b}=\left(\sqrt{3}; 1; 0\right)$. Góc giữa hai vectơ $\vec{a},$ $ \vec{b}$~là $n^{\circ}$ với $n$~là số thực. Giá trị của $n$~là bao nhiêu?
	 \shortans{120}
	 \loigiai{Ta có
		$$ 
		\cos \left(\vec{a}, \vec{b}\right)=\dfrac{0 \cdot \sqrt{3}+(-1) \cdot 1+0 \cdot 0}{\sqrt{0^{2}+(-1)^{2}+0^{2}} \cdot \sqrt{\left(\sqrt{3}\right)^{2}+1^{2}+0^{2}}}=\dfrac{-1}{2} \Rightarrow\left(\vec{a}, \vec{b}\right)=120^{\circ}. 
	$$} 
\end{bt}

% nguồn: Lop_12_Hinh_hoc_Chuong_2
\begin{bt}%[2H2H2-4]%[ma:0cf25985]
	Trong không gian $Oxyz$, cho $3$ điểm $A(2;0;0)$, $B(0;3;1)$; $C(-3;6;4)$. Gọi $M$ là điểm nằm trên đoạn $BC$ sao cho $MC=2MB$. Tính độ dài đoạn $AM$ (Kết quả được làm tròn ở chữ  số thập phân thứ nhất).
	\shortans{5,4} 
	\loigiai{ 
		Gọi tọa độ điểm  $M(a;b;c)$, do $M$ thuộc đoạn $BC$ và $MC=2MB\Leftrightarrow \overrightarrow{CM}=2\overrightarrow{MB}$ $(*)$.\\
		 Ta có $\overrightarrow{CM}=(a+3;b-6;c-4)$ và $\overrightarrow{MB}=(-a;3-b;1-c)$. \\
		 Do đó \\
		$(*)\Leftrightarrow \heva{&a+3=-2a\\&b-6=6-2b\\&c-4=2-2c}\Leftrightarrow \heva{&a=-1\\&b=4\\&c=2}\Rightarrow  M(-1;4;2)\Rightarrow\overrightarrow{AM}=(-3;4;2)$\\
		$\Rightarrow  AM=\sqrt{29}\approx 5{,}4$.
	 }
\end{bt}

% nguồn: Lop_12_Hinh_hoc_Chuong_2
\begin{bt}%[2H2H2-4]%[ma:636e9a31]
	Trong không gian  $Oxyz$, cho $\overrightarrow{a}=(3;1;-2)$ và $\overrightarrow{b}=(-2;0;-3)$. Tính tích vô hướng $\overrightarrow{a} \cdot(2 \overrightarrow{a}+\overrightarrow{b})$.
	\shortans{28} 
	\loigiai{ 
		 Ta có $\overrightarrow{a}=(3;1;-2) \Rightarrow 2 \overrightarrow{a}=(6;2;-4) \Rightarrow 2 \overrightarrow{a}+\overrightarrow{b}=(4;2;-7)$.\\
		Do đó $\overrightarrow{a} \cdot(2  \overrightarrow{a}+\overrightarrow{b})=12+2+14=28$.
	} 
\end{bt}

% nguồn: Lop_12_Hinh_hoc_Chuong_2
\begin{bt}%[2H2H2-4]%[ma:0cf3ecb6]
	Trong không gian $Oxyz$, cho hai vec tơ $\overrightarrow{u}=(1;1;0)$ và  $\overrightarrow{v}=(2;0;-1)$. Tính độ dài $\left|\overrightarrow{u}+2 \overrightarrow{v}\right|$ (\textit{kết quả làm tròn đến hàng phần trăm}).
	\shortans{5,48} 
	 \loigiai{
		Ta có $\left|\overrightarrow{u}\right|=\sqrt{2}$; $\left|\overrightarrow{v}\right|=\sqrt{5}$; $\overrightarrow{u} \cdot \overrightarrow{v}=1\cdot 2+1\cdot 0+0 \cdot  (-1)=2$.\\
		Suy ra $\left|\overrightarrow{u}+2 \overrightarrow{v}\right|^2=\left|\overrightarrow{u}\right|^2+4 \overrightarrow{u}\cdot  \overrightarrow{v}+4\left|\overrightarrow{v}\right|^2=2+4\cdot 2+4\cdot 5=30$.\\
		Vậy $\left|\overrightarrow{u}+2  \overrightarrow{v}\right|=\sqrt{30}\approx 5{,}48$.
	 }
\end{bt}

% nguồn: Lop_12_Hinh_hoc_Chuong_2
\begin{bt}%[2H2H2-4]%[ma:ff1d94a4]
	Trong  không gian với hệ tọa độ $Oxyz$, cho véc tơ $\overrightarrow{u}=(1;1;-2),\overrightarrow{v}=(1;0;m)$. Giá trị của $m$ (làm tròn đến hàng phần mười) để góc giữa $\vec{u}$, $\vec{v}$ bằng $45^{^\circ}$ là bao nhiêu?
	\shortans{-0,4} 
	\loigiai{ 
		Ta có 
		\allowdisplaybreaks 
		\begin{eqnarray*}
			 &&\left(\overrightarrow{u},\overrightarrow{v}\right)=45^\circ\\
			&\Leftrightarrow&\cos \left(\overrightarrow{u},\overrightarrow{v}\right)=\dfrac{\sqrt{2}}{2}\\ 
			&\Leftrightarrow&\dfrac{\overrightarrow{u}\cdot \overrightarrow{v}}{\left| \overrightarrow{u}\right|\cdot \left|  \overrightarrow{v}\right|}=\dfrac{\sqrt{2}}{2}\\
			 &\Leftrightarrow&\dfrac{1-2m}{\sqrt{6}\cdot \sqrt{1+m^2}}=\dfrac{1}{\sqrt{2}}\\
			&\Leftrightarrow&\heva{&1-2m\ge 0  \\&3m^2+3=1-4m+4m^2}\\
			&\Leftrightarrow&\heva{&m\le \dfrac{1}{2}  \\&m^2-4m-2=0}\\
			&\Leftrightarrow&m=2-\sqrt{6}\approx  -0{,}4494897.
		\end{eqnarray*}
		 Suy ra $m\approx -0{,}45$.
	} 
\end{bt}

% nguồn: Lop_12_Hinh_hoc_Chuong_2
\begin{bt}%[2H2H2-4]%[ma:79d85a90]
	Trong không gian $Oxyz$, cho $A(3;0;0)$, $B(0;3;0)$, $C(0;0;3)$. Điểm $M(a;b;c)$ trong không gian thỏa mãn $M$ không trùng với các điểm  $O$, $A$, $B$, $C$ và $\widehat{AMB}=\widehat{BMC}=\widehat{CMA}=90^\circ$. Khi đô tổng $a+b+c$ bằng
	\shortans{6} 
	\loigiai{ 
		Gọi $M(a;b;c)$ vì $\widehat{AMB}=\widehat{BMC}=\widehat{CMA}=90^\circ$ nên  $\heva{&\vv{AM}\cdot\vv{BM}=0\\ &\vv{CM}\cdot\vv{BM}=0\\&\vv{AM}\cdot\vv{CM}=0.} \, \, \left(*\right)$\\
		$\vv{AM}=(a-3;b;c)$, $\vv{BM}=(a;b-3;c)$, $\vv{CM}=(a;b;c-3)$.  \\
		 Từ biểu thức $(*)$ ta có biểu thức tọa độ
		\begin{eqnarray*}
			&\heva{&(a-3)a+b(b-3)+c^2=0\\&a^2+(b-3)b+c(c-3)=0\\&(a-3)a+b^2+(c-3)c=0}&\Leftrightarrow\heva{&a^2+b^2+c^2-3a-3b=0\, (1)\\&a^2+b^2+c^2-3b-3c=0\\&a^2+b^2+c^2-3a-3c=0}\\ 
			 &&\Leftrightarrow 3a+3b=3b+3c=3c+3a\\
			&&\Leftrightarrow  a=b=c.
		\end{eqnarray*}
		Thay vào biểu thức $(1)$  ta có $3a^2-6a=0\Leftrightarrow \hoac{&a=0\\&a=2} \Rightarrow\hoac{&M(0;0;0)\,(l)\\&M(2;2;2)\, (TM).}$\\
		Vậy  $a+b+c=6$.
	} 
\end{bt}

% nguồn: Lop_12_Hinh_hoc_Chuong_2
\begin{bt}%[2H2V2-4]%[ma:0695b324]
	Trong mặt phẳng tọa  độ $Oxyz$ cho các điểm $A(2;0;2)$, $B(0;2;0)$, $C(1;0;3)$. Gọi $M$ là điểm trong không gian thỏa mãn $M A^2+M C^2=M B^2$. Tính $M P$ với $P(3;-2;5)$. (\textit{Kết quả làm tròn đến hàng phần trăm})\shortans{4,90}
	\loigiai{Gọi $I(x;y;z)$ là điểm thỏa mãn $\overrightarrow{I  A}+\overrightarrow{I C}=\overrightarrow{I B}$\hfill $(*)$.\\
		Ta có $\overrightarrow{I A}=(2-x;-y;2-z)$; $\overrightarrow{I B}=(-x;2-y;-z)$;  $\overrightarrow{I C}=(1-x;-y;3-z)$.\\
		Khi đó $$(*)  \Leftrightarrow\heva{&2-x+1-x=-x \\& -y-y=2-y \\& 2-z+3-z=-z} \Leftrightarrow\heva{&x=3 \\& y=-2 \\& z=5} \Rightarrow I(3;-2;5) \equiv P.$$
		Suy  ra
		 \allowdisplaybreaks
		\begin{eqnarray*}
			\overrightarrow{I A}&=&(-1;2;-3) \Rightarrow I A^2=14;\\\overrightarrow{I B}&=&(-3;4;-5) \Rightarrow  I B^2=50;\\\overrightarrow{I C}&=&(-2;2;-2) \Rightarrow I C^2=12.
		\end{eqnarray*}
		Ta có $M A^2+M C^2=M B^2 \Leftrightarrow M  A^2+M C^2-M B^2=0$.\\
		Khi  đó
		\allowdisplaybreaks 
		\begin{eqnarray*}
			M A^2+M  C^2-M B^2&=&\left(\overrightarrow{M I}+\overrightarrow{I A}\right)^2+\left(\overrightarrow{M I}+\overrightarrow{I C}\right)^2-\left(\overrightarrow{M I}+\overrightarrow{I B}\right)^2\\
			&=&M I^2+I A^2+2 \overrightarrow{M I} \cdot  \overrightarrow{I A}+M I^2+I C^2+2 \overrightarrow{M I} \cdot \overrightarrow{I C}\\&&-M I^2-I B^2-2 \overrightarrow{M I} \cdot \overrightarrow{I B}
			\\&=&M I^2+\left(I A^2+I C^2-I B^2\right)+2 \overrightarrow{M I}\left(\overrightarrow{I A}+\overrightarrow{I  C}-\overrightarrow{I B}\right)=0\\
			 \Leftrightarrow M P^2+(14+12-50)&=&0 \Leftrightarrow M P^2=24 \Rightarrow M P=2 \sqrt{6}.
		\end{eqnarray*}
	 }
\end{bt}

% nguồn: Lop_12_Hinh_hoc_Chuong_2
\begin{bt}%[2H2V2-4]%[ma:fc7a26c1]
	Trong  không gian $Oxyz$, cho hai điểm $A(2;2;-2)$ và $B(-3;-3;3)$. Xét điểm $M$ thay đổi sao cho $3MA-2MB=0$. Giá trị lớn nhất của độ dài đoạn thẳng $OM$ (kết quả làm tròn đến hàng phần chục) bằng bao nhiêu?
	\shortans{20,8} 
	 \loigiai{
		Từ  giả thiết, ta có
		\begin{eqnarray*}
			&&  9 \overrightarrow{AM}^2=4 \overrightarrow{BM}^2 \\
			 &\Leftrightarrow &9(\overrightarrow{OM}-\overrightarrow{OA})^2=4(\overrightarrow{OM}-\overrightarrow{OB})^2\\
			&\Leftrightarrow &9\left(OM^2+OA^2-2 \overrightarrow{OM}  \cdot \overrightarrow{OA}\right)=4\left(OM^2+OB^2-2 \overrightarrow{OM} \cdot \overrightarrow{OB}\right)\\
			&\Leftrightarrow &5OM^2=4OB^2-9OA^2+2 \overrightarrow{OM}(9\overrightarrow{OA}-4\overrightarrow{OB}). \quad(*) 
		\end{eqnarray*}
		Khi đó, độ dài đoạn thẳng $OM$ lớn nhất khi và chỉ khi hai $\overrightarrow{OM}$ và $\overrightarrow{u}=9 \overrightarrow{OA}-4 \overrightarrow{OB}$ là các véc-tơ cùng hướng, với $\overrightarrow{u}=9 \overrightarrow{OA}-4 \overrightarrow{OB}=(30;30;-30)$.\\ 
		Đặt $\overrightarrow{OM}=t(1;1;-1)$ với $t \neq 0$, từ $\left(*\right)$ ta được $15t^2=0+180t \Rightarrow  t=\dfrac{180}{15}=12$.\\
		Vậy  $\overrightarrow{OM}=12(1;1;-1) \Rightarrow OM=12 \sqrt{3} \approx 20{,}8$.
	 }
\end{bt}

% nguồn: Lop_12_Hinh_hoc_Chuong_2
\begin{bt}%[2H2C2-4]%[ma:b21ec347]
	Trong không gian, cho hình chóp $ S. ABCD$ có đáy là hình bình hành tâm $O$, gọi $M$ là điểm thay đổi trên $SO$. Khi biểu thức $ P=M{S^2}+M{A^2}+M{B^2}+M{C^2}+M{D^2}$ đạt giá trị nhỏ nhất, tính tỉ  số $\dfrac{SM}{SO}$. (Kết quả làm tròn đến hàng phần chục)
	\shortans{0,8} 
	 \loigiai{
		Gọi  $I$ là điểm thỏa mãn $\vec{SI}=4\vec{IO}$.\\
		$\begin{aligned}
		 P= &\, (\vec{MI}+\vec{IS})^2+(\vec{MI}+\vec{IA})^2+(\vec{MI}+\vec{IB})^2+(\vec{MI}+\vec{IC})^2+(\vec{MI}+\vec{ID})^2\\
		=  &\, 5M{I^2}+I{S^2}+I{A^2}+I{B^2}+I{C^2}+I{D^2}+2\vec{MI}\left(\vec{IS}+\vec{IA}+\vec{IB}+\vec{IC}+\vec{ID}\right)\\
		= &\,5M{I^2}+I{S^2}+I{A^2}+I{B^2}+I{C^2}+I{D^2}+2\vec{MI}\left(\vec{IS}+4\vec{IO}+\vec{OA}+\vec{OB}+\vec{OC}+\vec{OD}\right) \\ 
		= &\,5M{I^2}+I{S^2}+I{A^2}+I{B^2}+I{C^2}+I{D^2}. \,(\text{do }  \vec{SI}=4\vec{IO} \text{ và } \vec{OA}+\vec{OB}+\vec{OC}+\vec{OD}=\vec{0})
	\end{aligned}$
	Vậy $P_{\text{ min}}\Leftrightarrow M\equiv  I\Rightarrow\dfrac{SM}{SO}=\dfrac{4}{5}=0{,}8$.
} 
\end{bt}


%%%%%%%%%%%%%%%% TỰ LUẬN %%%%%%%%%%%%%%%%

% nguồn: Lop_12_Hinh_hoc_Chuong_2
\begin{bt}%[2H2H2-4]%[ma:87426a30]
	Trong không gian $Oxyz$, cho ba vectơ $\overrightarrow{a}=(3;1;2)$, $\overrightarrow{b}=(-3;0;4)$ và $\overrightarrow{c}=(6;-1;0)$. 
	\begin{enumerate}
		\item Tìm tọa độ của các vectơ  $\overrightarrow{a}+\overrightarrow{b}+\overrightarrow{c}$ và $2 \overrightarrow{a}-3 \overrightarrow{b}-5 \overrightarrow{c}$.
		\item Tính các tích vô hướng  $\overrightarrow{a} \cdot(-\overrightarrow{b})$ và $(2 \overrightarrow{a}) \cdot \overrightarrow{c}$.
	\end{enumerate}
	\loigiai{ 
		\begin{enumerate}
			\item Tọa độ của  vectơ $\overrightarrow{a}+\overrightarrow{b}+\overrightarrow{c}$ là $(3-3+6;1+0-1;2+4+0)=(6;0;6)$.\\
			Ta có $2 \overrightarrow{a}=(6;2;4)$,  $3\overrightarrow{b}=(-9;0;12)$, $5\overrightarrow{c}=(30;-5;0)$.\\
			Tọa  độ của vectơ $2 \overrightarrow{a}-3 \overrightarrow{b}-5 \overrightarrow{c}$ là $$(6+9-30;2+5;4-12)=(-15;7;-8).$$
			 \item Ta có $-\overrightarrow{b}=(3;0;-4)$.\\
			Do đó $\overrightarrow{a} \cdot(-\overrightarrow{b})=3\cdot 3+1\cdot 0+2\cdot(-4)=1$ và $\left(2 \overrightarrow{a}\right)\cdot \overrightarrow{c}=6 \cdot  6+2 \cdot(-1)+4 \cdot 0=34$.
		\end{enumerate}
	} 
\end{bt}

% nguồn: Lop_12_Hinh_hoc_Chuong_2
\begin{bt}%[2H2H2-4]%[ma:ed62bbbc]
	 \immini[thm]{
	Cho hình chóp $S.ABCD$ có đáy $ABCD$ là hình chữ nhật và $SA$ vuông góc với $(ABCD)$. Giả sử  $SA=2$, $AB=3$, $AD=4$. Xét hệ tọa độ $Oxyz$ với $O$ trùng $A$ và các tia $Ox$, $Oy$, $Oz$ lần lượt trùng với các tia $AB$, $AD$, $AS$.
	\begin{enumerate}
		\item Xác định tọa  độ của các điểm $S$, $A$, $B$, $C$, $D$.
		\item Tính $BD$ và  $SC$.
		\item Tính  $\left(\overrightarrow{BD},\overrightarrow{SC}\right)$.
	\end{enumerate}
	} 
	{ 
	\begin{tikzpicture}[scale=1,font=\footnotesize,line join=round,
				 line cap=round,>=stealth]
			\path (0:0) coordinate (A)--(3,0) coordinate (D)--++(-140:1) coordinate  (C)--++(180:3) coordinate (B)
				 (A)--++(90:3.2) coordinate (S); \draw[->](S)--++(90:.8) coordinate (z) node[right]{$z$};
			\draw[->](D)--++(0:1) coordinate  (y) node[above]{$y$};
			\draw[->](B)--++(-140:.8) coordinate (x)  node[right]{$x$};
			 \draw[dashed] (B)--(D)
				 (S)--(A)--(D)
				 (A)--(B);
			\draw (S)--(B)--(C)--(D)--cycle  (S)--(C);
			\draw (A) node[above  right]{$O$};
			\foreach \x/\g in {A/150,B/150,C/0,D/30,S/180} 
			\fill (\x) circle (1pt) 
				($(\x)+(\g:3mm)$) node{$\x$}; 
		\end{tikzpicture}
	} 
	\loigiai{ 
		\begin{enumerate}
			\item Vì $A$ trùng gốc tọa độ nên $A(0;0;0)$, $B\in Ox$ và  $AB=3$ nên $B(3;0;0)$.\\
			Vì $D\in Oy$ và $AD=4$ nên $D(0;4;0)$, $S\in Oz$ và $AS=2$  nên $S(0;0;2)$.\\
			Vì hình chiếu của $C$ lên các trục $Ox$, $Oy$, $Oz$ lần lượt là  $B$, $D$, $A$ nên $C(3;4;0)$.
			\item Ta  có $\overrightarrow{BD}=(0-3;4-0;0-0)=(-3;4;0)$, suy ra $BD=\left|\overrightarrow{BD}\right|=\sqrt{(-3)^2+4^2+0^2}=5$.\\
			 Ta có $\overrightarrow{SC}=(3-0;4-0;0-2)=(3;4;-2)$, suy ra $SC=\left|\overrightarrow{S C}\right| =\sqrt{3^2+4^2+(-2)^2}=\sqrt{29}$.
			\item Ta có $\cos \left(\overrightarrow{BD}, \overrightarrow{SC}\right) =\dfrac{\overrightarrow{BD} \cdot \overrightarrow{SC}}{\left| \overrightarrow{BD}\right|\cdot\left| \overrightarrow{SC}\right|}=\dfrac{(-3) \cdot  3+4\cdot 4+0 \cdot(-2)}{5 \sqrt{29}}=\dfrac{7}{5 \sqrt{29}}$.\\
			Suy ra $\left(\overrightarrow{BD},\overrightarrow{SC}\right) \approx 74{,}9^{\circ}$. 
		\end{enumerate}
	} 
\end{bt}
